% Transcription — fonds Alexandre Grothendieck, shelfmark 9, batch 1 (pages 1-20).
% Established from the facsimile archives/batches/9/batch-01.pdf, cut from a
% copy of 9.pdf that was fetched through the project's own relay
% (https://grothendieck-relay.onrender.com/source/9.pdf) rather than from
% grothendieck.umontpellier.fr directly; per the skill transcribe-grothendieck.
% The batch file carries no cover sheet: PDF page k is the archivists' page k,
% and their pencil numbers bottom-left agree wherever legible. The folder is
% twenty-eight pages; this is the first of two batches (pages 21-28 are
% batch 2, transcribed in parallel by another pass; its notation was read
% and is followed here).
%
% Pass: Opus 5.5 (claude-opus-5-5), 2026-09-26 — first pass, unchecked against the pages by a human.
%
% The folder sits in the inventory group "4-9" « Cristaux », dated
% [à partir de 1958- vers 1971]; the folder has its own date, which
% \dating{} carries verbatim, as batch 2 does.
%
% Pages skipped: 3, a sheet of Pisa letterhead
% (Università di Pisa, Istituto Matematico « Leonida Tonelli ») carrying
% nothing but show-through; 5, a leaf of the same letterhead with a list of
% Italian phrases and a street sketch, no mathematics. Page 8 was scanned
% upside down and is read turned.
%
% Covers (no \page marker, his words become the \section heading, with a
% \note): page 1 (buff card), « Dualité des modules de Dieudonné de gpes de
% B.T. duals l'un de l'autre ..... Le module de Dieudonné réduit M ⊗_W k »;
% page 9 (pink card), « Caractères des groupes de B.T. à réduction
% semi-stable virtuelle essentielle »; page 14 (buff card), a struck
% « Déformations des groupes de B.T. ordinaires », then « Détermination des
% "stratifications canoniques" sur les déformations » and « L'isom. de nature
% globale » with the boxed det T_p(A)_ét ≃ det T_p(A*)_ét.
%
% Pages summarised: none. Pages 10-13 are a typescript (typed in blue, no
% typed pagination, not identified as published), with his ink title and
% corrections; it is given in full, his ink interventions flagged in notes.
%
% What the batch is about. Four unpaginated runs.
%   2       (Pisa letterhead) G p-divisible over k of char. p, dual G*; the
%           two exact sequences 0 → t_G^∨ → H(G) → t_{G*} → 0 and
%           0 → t_{G*}^∨ → t_{G~} → t_G → 0 (G~ the universal extension by
%           V(t_{G*})), with F_0, V_0 between H(G)^{(p)} and H(G); Ker F_0 =
%           Im V_0 = ω_G^{(p)}; the Dieudonné comparison; a trichotomy of
%           nilpotence of F_0, V_0 against G, G* infinitesimal.
%   4, 6, 7 The Dieudonné ring A, A_n = A/(F^{n+1}, V^{n+1}): the trace form
%           φ_n : A_n → W_n (φ_n(1) = 1, φ_n(F^i) = φ_n(V^i) = 0), the
%           bimodule I = lim I_n, the dualising functors D and
%           D' = Hom_{W_n}(-, W_n), the pairing <x, y> = φ_n(xy) and its
%           rules under F, V, λ.
%   8       (cancelled by three strokes) Question: how to define an
%           integrable connection on R f_*(Ω^•_{X/S}) directly; a lemma on a
%           G-torsor P with θ : P → Z(Hom^•(K,L)^0) giving φ : K → L in D(X).
%   10-13   The typescript « Caractère d'un groupe de B.T. sur un corps
%           local »: monodromy condition, field of virtual good reduction,
%           Grothendieck groups K_BT(K;K'), homomorphisms (1)-(5) to
%           R_{W(k')}(I)^{G/I} and Cent(I, W(k'))^G, the character of M;
%           comparison with the modular character of M(0) = _pM on
%           p-regular elements of inertia, the counterexample Z_p(1), the
%           « cas test » of two elliptic curves with isomorphic _pE, and the
%           elliptic case (his pencil: « faux par Tate »).
%   15-20   Stratifications on DR^1 for deformations of an abelian variety
%           in char. p > 0: the question, compatibility with F, V, the
%           conclusion that A_0 must be ordinary, the stable decomposition
%           DR^1 ≃ t_A^∨ + t_{A*} for A ordinary, t_A ≃ L_A ⊗_{F_p} O_S with
%           L_A = (_pA*)_ét^∨; a (partly cancelled) N.B. on a polarisation
%           as a section of Λ^2 DR^1 and its middle component, the tangent
%           map of φ~; « Re N.B. »: ω_A = det(t_A)^{-1} and the section
%           σ_A ∈ Γ ω_A^{(p-1)} from t_A^{(p)} → t_A — the Hasse invariant,
%           whose naming opens batch 2 on page 21.
%
% Notation fixed once for the batch, following batch 2: A* the dual;
% \underline{t}_A, \underline{\mathcal{O}}_S underlined as written (the
% underline is not always visible on p. 2 but is kept); ω_A; DR^1(A/S);
% (_pA)_ét with {}_{p}A for his left subscript; \tilde\varphi : A → A* for
% the map of a polarisation; σ_A. His curly L of pp. 18-19 is \mathcal{L};
% his script H of pp. 15-16 \mathcal{H}; H^{10}, H^{01} for the Hodge pieces.
% The typescript's hand-inked Greek (ν, μ, ℓ, φ) is in math mode.
%
% Readings to check first: all the prose of pp. 2, 15, 16, 17 (fast
% drafting: formulas carry, connective prose largely does not); the second
% word of the heading on p. 2 (« Cas ... »); the right-hand exponent labels
% on p. 7; the struck parenthesis and the cancelled block on p. 20; on p. 19
% the extent of the diagonal cancellation; « (quitte : » on p. 10.

\documentclass[11pt,a4paper]{article}
% Le préambule commun à toutes les transcriptions du fonds Grothendieck.
%
% Il tient en une idée : **l'incertitude se compose, elle ne se résout pas.**
% Une transcription qui lisse un mot douteux a détruit la seule chose pour
% laquelle on la faisait — un lecteur doit pouvoir savoir, à chaque endroit,
% ce qui était sur le feuillet et ce qui a été ajouté. Les six macros qui
% suivent sont donc l'appareil critique, et non de la décoration.
%
% Les mêmes commandes sont reconnues par `scripts/render.mjs`, qui produit la
% vue de lecture du site. Une macro ajoutée ici doit l'être aussi là-bas,
% faute de quoi le rendu échouera bruyamment — ce qui est voulu : un rendu
% silencieusement faux est pire qu'un rendu absent.

\ProvidesPackage{grothendieck}[2026/08/08 Transcriptions du fonds Grothendieck]

\usepackage{amsmath,amssymb,amsthm}
\usepackage{mathrsfs}
\usepackage{stmaryrd}
% Les diagrammes commutatifs sont partout dans ce fonds : c'est la langue
% même de la géométrie algébrique de Grothendieck, et une transcription qui
% les rendrait en prose ne transcrirait rien.
\usepackage{tikz-cd}
% Français par défaut — c'est la langue des feuillets — mais l'anglais est
% chargé aussi : la traduction et le résumé sont des documents anglais, et la
% typographie française (espaces avant les deux-points) y serait une faute.
% Ces documents basculent par \selectlanguage{english} en tête de corps.
\usepackage[english,french]{babel}
\usepackage{fontspec}
\usepackage{geometry}
\geometry{a4paper,margin=25mm,marginparwidth=28mm}
\usepackage{marginnote}
\usepackage{xcolor}
% ulem, not soul. `soul`'s \st and \ul are built on a character-by-character
% scan that XeLaTeX's font machinery breaks: the document fails with an
% undefined control sequence rather than degrading. ulem does the same two jobs
% and is Unicode-engine safe. `normalem` stops it hijacking \emph, which the
% transcriptions use for Grothendieck's own emphasis.
\usepackage[normalem]{ulem}
\usepackage{hyperref}
\hypersetup{colorlinks=true,linkcolor=black,urlcolor=black,pdfborder={0 0 0}}

% --- Métadonnées du lot -----------------------------------------------------
% Reprises telles quelles par le script de rendu, qui les affiche en tête de
% la vue de lecture. Elles fixent ce que le fichier prétend couvrir : sans
% elles, une transcription flotte sans point d'attache dans un fonds de seize
% mille feuillets.

\newcommand{\folder}[1]{\gdef\grothFolder{#1}}
\newcommand{\batch}[1]{\gdef\grothBatch{#1}}
\newcommand{\leaves}[2]{\gdef\grothFirst{#1}\gdef\grothLast{#2}}
\newcommand{\foldertitle}[1]{\gdef\grothFoldertitle{#1}}
\gdef\grothFolder{?}\gdef\grothBatch{?}\gdef\grothFirst{?}\gdef\grothLast{?}\gdef\grothFoldertitle{}

% --- Le résumé --------------------------------------------------------------
%
% Il ouvre la lecture modernisée, sous le titre « Résumé », et il est composé
% en petit. Ce n'est pas un ornement : c'est le seul passage du document qui
% ne soit pas des mathématiques, et un lecteur doit pouvoir le reconnaître
% avant de l'avoir lu — puis le sauter s'il connaît déjà le sujet, ou s'y
% arrêter s'il ne le connaît pas. Le filet qui le suit marque la reprise du
% corps.

\newenvironment{resume}
  {\small\setlength{\parskip}{0.45em}}
  {\par\medskip\hrule\medskip}

% --- Mots-clés ---------------------------------------------------------------
%
% En anglais, en clôture du résumé de la lecture modernisée : le vocabulaire
% moderne sous lequel chercher ce que les feuillets construisent. Le manifeste
% du site (`npm run manifest`) les extrait pour étiqueter la cote entière —
% cette ligne est la source unique des tags, il n'y a pas de fichier à côté.
\newcommand{\keywords}[1]{\par\smallskip\noindent
  {\footnotesize\textsc{Keywords} --- \textit{#1}}\par}

% --- L'appareil critique ----------------------------------------------------

% Le numéro de feuillet, en marge. C'est l'ancre qui permet de régler un
% désaccord sur une formule en regardant *un* feuillet plutôt qu'une chemise
% de six cents. Le site s'en sert aussi pour tourner les pages du fac-similé
% quand on fait défiler la transcription.
\newcommand{\leaf}[1]{%
  \marginnote{\footnotesize\color{gray!70}\textsf{#1}}%
  \label{leaf:#1}\ignorespaces}

% Illisible : on ne devine pas. Un mot inventé qui se lit comme les autres est
% le pire résultat possible.
\newcommand{\ill}{\textcolor{red!60!black}{[\,\ldots\,]}}

% Lecture douteuse : le mot est donné, et signalé comme incertain.
\newcommand{\uncertain}[1]{\uline{#1}}

% Ajout de l'éditeur — une lettre manquante, un mot sous-entendu.
\newcommand{\add}[1]{\textcolor{blue!50!black}{[#1]}}

% Biffé par Grothendieck lui-même. Ce qu'il a rayé fait partie du document :
% une rature dit souvent où la pensée a changé de direction.
\newcommand{\struck}[1]{\sout{#1}}

% Note du transcripteur, sur l'état matériel du feuillet.
\newcommand{\note}[1]{\par\smallskip{\small\color{gray!60!black}\textit{[#1]}}\par\smallskip}

% Note marginale de Grothendieck, distincte du corps du texte.
\newcommand{\marginal}[1]{\par\smallskip\noindent\hspace{1em}%
  {\small\color{gray!60!black}#1}\par\smallskip}

% --- Titre ------------------------------------------------------------------

\AtBeginDocument{%
  \begin{center}
    {\large\bfseries Fonds Alexandre Grothendieck --- cote n\textsuperscript{o}~\grothFolder}\\[2pt]
    {\itshape\grothFoldertitle}\\[4pt]
    {\small feuillets \grothFirst--\grothLast\ (lot \grothBatch)}\\[2pt]
    {\footnotesize\color{gray!60!black}Transcription établie d'après le fac-similé numérisé
     par l'Université de Montpellier.\\
     Les lectures incertaines sont soulignées, les illisibles notées [\,\ldots\,],
     les ajouts entre crochets.}
  \end{center}
  \vspace{1em}}

\endinput


\folder{9}
\batch{1}
\pages{1}{20}
\dating{[vers 1969]}
\watermark{Édition de démonstration}
\foldertitle{Groupe de Barsotti-Tate et cristaux (jusqu’en été 1969) : notes manuscrites (s.d.)}

\begin{document}

\section*{Dualité des modules de Dieudonné de gpes de B.T.\ duals l'un de l'autre}

\note{inscrit à l'encre en haut à droite du feuillet de garde (page 1), avec
en dessous « Le module de Dieudonné réduit $M \otimes_W k$ » ; le feuillet ne
reçoit pas de numéro de page}

\page{2}

$k$ corps de car.\ $p > 0$

$G$ gpe $p$-divisible, $G^*$ dual

$g = \dim G$, $g^* = \dim G^*$

\[
\begin{array}{ccccccccc}
0 & \to & \underline{t}_G^{\vee} & \to & H(G) & \to & \underline{t}_{G^*} & \to & 0 \\
 & & \| & & {\scriptstyle H(V) = V_0}\downarrow\;\uparrow{\scriptstyle H(F) = F_0} & & & & \\
 & & \omega_G & & & & & & \\
0 & \to & \underline{t}_G^{\vee(p)} & \to & H(G)^{(p)} & \to & \underline{t}_{G^*}^{(p)} & \to & 0 \\
 & & \| & & & & & & \\
 & & \omega_G^{(p)} & & & & & & \\
 & & \dim g & & \dim g + g^* & & \dim g^* & &
\end{array}
\]
\note{les deux suites sont écrites l'une sous l'autre, $V_0$ descendant de
$H(G)$ vers $H(G)^{(p)}$, $F_0$ remontant ; « $\dim g$ », « $\dim g+g^*$ »,
« $\dim g^*$ » sont inscrits sous les trois termes, tels qu'écrits (pour
$\dim = g$, etc.) ; le « $H(V)$ » est surchargé. On rend $t$ souligné comme
en page~21 (lot~2), bien que le soulignement ne soit pas toujours visible ici}

$\tilde G$ \uncertain{groupe} \ill{}\note{deux ou trois mots barrés d'un
trait, sur l'en-tête en transparence}
\[
0 \to \mathbb{V}(\underline{t}_{G^*}) \to \tilde G \to G \to 0
\]

\marginal{dual de}
\[
\begin{array}{ccccccccc}
0 & \to & \underline{t}_{G^*}^{\vee} & \to & \underline{t}_{\tilde G} & \to & \underline{t}_G & \to & 0 \\
 & & \| & & {\scriptstyle V_0}\downarrow\;\uparrow{\scriptstyle F_0} & & & & \\
 & & \omega_{G^*} & & & & & & \\
0 & \to & \underline{t}_{G^*}^{\vee(p)} & \to & \underline{t}_{\tilde G}^{(p)} & \to & \underline{t}_G^{(p)} & \to & 0 \\
 & & \| & & & & & & \\
 & & \omega_{G^*}^{(p)} & & & & & &
\end{array}
\]
\note{les deux diagrammes sont côte à côte, séparés par un trait vertical,
« dual de » écrit au-dessus du trait}

[\textbf{NB} $H(G) = \underline{t}_{\tilde G^*}$]\note{l'indice est
recouvert d'une tache d'encre ; seul « $G^*$ » en dessous se lit, et la
lecture $\tilde G^*$ s'appuie sur la même égalité, lisible, plus bas}

(1) suite exacte pour $\underline{t}_{\tilde G}$ et
$\underline{t}_{\tilde G^*}$ \uncertain{transposée} l'une de
\struck{\ill{}} l'autre, \ill{} \ill{} \ill{} \uncertain{fonctoriel},
\ill{} \uncertain{laquelle} $F_0$ et $V_0$ \ill{} \ill{} \ill{}
\ill{} $p$) \uncertain{s'interchangent}.

(2) \struck{Donc} La suite exacte \ill{} \ill{},
\[
\begin{aligned}
\operatorname{Ker} F_0 &= \operatorname{Im} V_0 = \underline{t}_G^{\vee(p)} = \omega_G^{(p)} \\
\operatorname{Ker} V_0 &= \operatorname{Im} F_0 = \text{ce que ça \ill{}}
\end{aligned}
\]
(\ill{} \ill{}) \ill{} \ill{} \ill{}\note{la fin de la ligne, à droite
de « $\operatorname{Im} F_0 =$ », est une parenthèse de prose rapide que
l'on n'a pas su lire}

(3) $F_0$ et \struck{$V$} $V_0$ \ill{} sur $H(G)$

? Dieudonné : \struck{\ill{}} $G$ réduit mod $p$.
\[
\underline{t}_{G^*}^{(p)} \xrightarrow{F_0} H(G) \to \underline{t}_{G^*}
\]
\uncertain{correspond} : le \ill{} \uncertain{dérivation} \ill{}
[et correspond : Frobenius dans $H^2(G, \mathbb{G}_m)$] \ill{} dans
$H(G) = \underline{t}_{\tilde G^*}$\note{deux lignes serrées à droite de la
suite, avec un crochet ; ordre et coupure des deux lignes incertains}

(4) Le \uncertain{comparaison} \struck{\ill{}}

\emph{\uncertain{Cas} \ill{}}\note{deux mots soulignés en tête des trois lignes qui suivent ; le second ne se lit pas}
\begin{itemize}
\item $G$ \uncertain{purement} \uncertain{infinitésimal}, i.e.\ $G^*$
\uncertain{sans tore} $\iff$ $F_0$ \uncertain{nilpotent} ;
\item $G^*$ \uncertain{purement} \uncertain{infinitésimal}, i.e.\ $G$ sans
\uncertain{composante} \uncertain{torique} $\iff$ $V_0$ \uncertain{nilpotent} ;
\item $G$ et $G^*$ \uncertain{purement} \uncertain{infinitésimaux} $\iff$
$F_0$ et $V_0$ \uncertain{nilpotents}.
\end{itemize}

\page{4}

\[
\langle x, y \rangle = \varphi_n(xy) \qquad \varphi_n : A_n \to W_n
\]
\begin{itemize}
\item[a)] $\varphi_n$ est \struck{\ill{}} linéaire : \uncertain{gauche} et
\uncertain{droite} \uncertain{sur} $W_n$\note{les deux mots soulignés, le
premier deux fois}
\item[b)] $\varphi_n(Fz) = \struck{\ill{}}\, \varphi_n(zF)^{\pi}$
\item[c)] $\varphi_n(zV) = \struck{\ill{}}\, \varphi_n(Vz)^{\pi}$
\end{itemize}
\marginal{encadré, en haut à droite :
$\varphi_n(1) = 1$, $\varphi_n(F^i) = 0$, $\varphi_n(V^i) = 0$ $(i \geqslant 1)$ ;
en dessous : $\exists !$ $\varphi_n$}

d) \uncertain{Les} $\varphi_n$ \uncertain{se} \ill{} ($\varphi_n$ \ill{}
\uncertain{variable})

\struck{$F^iV^i \mapsto a$}

\struck{$\varphi\bigl(\sum_{0 \leqslant i, j \leqslant n} a_{ij} F^i V^j\bigr)
= \ill{}$}\note{sous la somme une accolade marquée $z$}

\[
\varphi_n\Bigl(a + \sum_{1}^{\to} \lambda_i F^i + \sum_{1}^{\to}
\struck{\ill{}}\, V^i \mu_i\Bigr) = \varphi_n(z) = a
\]
\note{l'argument est marqué $z$ par une accolade. Autour, un premier essai
annulé par des traits en zigzag : $\varphi_n(1) = a$, $\sum_0 \lambda_i
F^i$, $\varphi_n(F\ldots$, $\varphi_n(Fz) = p\ldots$, $\varphi_n(zF) =$ ;
on ne le détaille pas}

\[
\begin{aligned}
\varphi_n(\lambda z) &= \lambda \varphi_n(z) \qquad \text{OK} \\
\varphi_n(z\lambda) &= \varphi_n\Bigl(\sum \lambda_{ij} \lambda^{\pi^{i-j}}
F^i V^j\Bigr) = \varphi_n(z)\lambda \qquad \text{OK.}
\end{aligned}
\]
\struck{$\varphi_n(Fz) = \varphi\bigl(\sum \lambda_{ij}^{\pi} F^{i+1}V^j\bigr)
= \ill{}$} / $\varphi_n(zF) = \varphi_n\bigl(\sum \lambda_{ij} F^{i+1} V^j\bigr)$

$\varphi_n(Fz) = \varphi_n(zF) =$\note{la page s'arrête sur ce signe égal}

\page{6}

\[
A_n = A/(F^{n+1}, V^{n+1})
\]
$I_n = A_n$ \uncertain{regardé} comme \emph{\uncertain{bimodule}} sur $A$
\[
I = \varinjlim_n I_n \qquad I_n \xrightarrow{\ p\ } I_{n+1}
\]
\note{les exposants $n+1$ de $A/(F^{n+1}, V^{n+1})$ sont lus sur un trait
rapide ; un trait horizontal sépare ensuite ce bloc de la suite}

$D : M \mapsto \operatorname{Hom}(M, I)$ \struck{\ill{}} foncteur
\uncertain{dualisant}.

\emph{Attention}, il \uncertain{transforme} un module à \uncertain{gauche}
\uncertain{en} \uncertain{module} à droite, et vice versa. Mais on
\uncertain{peut} utiliser l'\uncertain{involution} $s$ de $A$ échangeant $F$
et $V$, pour passer de l'un à l'autre.

\struck{$D(M)$} $A_n$ est fini sur $W_n$. Un autre foncteur dualisant est
obtenu en faisant (si $M$ tué par $F^n, V^n$)
\[
D' : M \mapsto \operatorname{Hom}_{W_n\text{-}\mathrm{mod.}}(M, W_n)
\]
\note{sous « $\operatorname{Hom}$ » un indice en toutes lettres,
« $W_n$-modules », lecture incertaine}
et définissons sur $D'(M)$ la structure de $A$-module à droite par
« \uncertain{transposition} ».

\struck{$D'(M) = $} $\operatorname{Hom}(A_n, W_n)$ structure de
\uncertain{bimodule}

\struck{$\operatorname{Hom}_{W_n}(M, W_n) = \operatorname{Hom}_{A_n}\bigl(M,
\operatorname{Hom}_{W_n}(A_n, W_n)\bigr)$}\note{la ligne est barrée ; sous le
dernier terme, une accolade renvoie à « structure de bimodule naturelle », et
un long trait courbe ramène le tout à la dernière ligne}

Je dis qu'on a un \uncertain{isom.}
\[
I_n = A_n \simeq \operatorname{Hom}_{W_n}(A_n, W_n)
\]
\struck{$u F = u \circ F$}

\page{7}

$D'(M) = \operatorname{Hom}_{W_n}(M, W_n)$ \uncertain{module} à
\emph{droite} sur $A_n$

\struck{$uF\ \ Fu\ \ (uF)(x) = (u(F(x)))^{\pi}$}

\[
\begin{aligned}
u \cdot F &= \pi^{-1}(u \circ F) \\
u \cdot V &= \pi(u \circ F) \\
u \cdot \lambda &= u \circ \lambda
\end{aligned}
\qquad\Bigg|\qquad
\begin{aligned}
u \cdot F \cdot V &= \pi\pi^{-1}\, u \circ F \circ V = p\,u \\
u\,VF &= p\,u
\end{aligned}
\]
\note{le second membre de $u \cdot V$ porte bien « $u \circ F$ », tel
qu'écrit}

\emph{\uncertain{bimodule}} sur $A_n$

\[
D'(A_n) = \operatorname{Hom}_{W_n}\bigl((A_n)_g, W_n\bigr)
\]
\[
\begin{aligned}
u \cdot F &= \pi^{-1}(u \circ F_g) \\
u \cdot V &= \pi(u \circ V_g) \\
u \cdot \lambda &= u \circ \lambda_g \\
Fu &= u \circ F_d \\
Vu &= u \circ F_g \\
\lambda u &= u \circ \lambda_d
\end{aligned}
\qquad
u\lambda = \lambda \cdot u = u \circ \lambda = \lambda \circ u
\]
\note{les indices $g$, $d$ : action à gauche, à droite. « $Vu = u \circ
F_g$ » tel qu'écrit ; le $V$ surcharge une autre lettre}

\emph{\uncertain{Attention}} les deux structures de $W_n$-module sur
$I$ \ill{} \ill{} \ill{} \ill{} \ill{} \ill{}\note{deux lignes de prose
rapide que l'on n'a pas su lire}

\[
A_n \overset{?}{\xrightarrow{\ \sim\ }} \operatorname{Hom}_{W_n}\bigl((A_n)_g, W_n\bigr)
\qquad \text{\uncertain{isom.}\ de \uncertain{bimodules}}
\]

\[
\boxed{A_n \times A_n \longrightarrow W_n}
\qquad
(x, y) \mapsto \langle x, y \rangle
\qquad \text{\uncertain{bimodule}}
\]
\[
\begin{aligned}
&\langle \lambda x, y \rangle = \lambda \langle x, y \rangle & \lambda &\in W_n \\
&\langle x, y\lambda \rangle = \langle \lambda x, y \rangle & \lambda &\in W_n \\
&\langle x, yF \rangle = \pi^{-1} \langle Fx, y \rangle \\
&\langle x, yV \rangle = \pi \langle Vx, y \rangle \\
&\langle x, \lambda y \rangle = \langle x\lambda, y \rangle \\
&\langle x, Fy \rangle = \langle xF, y \rangle \\
&\langle x, Vy \rangle = \langle xV, y \rangle
\end{aligned}
\]
\note{marques « $x$ », « $x$ », puis « $\lambda$ », « $\alpha$ », « $x$ »
dans la marge gauche, en face des lignes ; sous la dernière, un « $F$ »
biffé}

\marginal{à droite, deux accolades : la première, sur les lignes 2 à 4,
« \uncertain{assure} \uncertain{que} $A_n \to \operatorname{Hom}_{W_n}((A_n)_g,
W_n)$ \uncertain{est} $A_n$-\uncertain{linéaire} \emph{à droite} » ; la
seconde, sur les trois dernières, « \uncertain{assure} \ill{} \ill{}
$A_n$-\uncertain{lin.} », suivi de \struck{$\langle x, y \rangle$}
\struck{$\langle x, y \rangle = \langle 1, xy \rangle$} i.e.\ $\langle x, y
\rangle = \langle 1, xy \rangle = \langle xy, 1 \rangle$}

\[
A_n \underset{A_n}{\otimes} A_n \longrightarrow W_n \qquad
W_n\text{-bilinéaire}
\]
\note{le tout encadré ; sous le produit tensoriel et sous $W_n$, deux
renvois « $W_n$ » et un mot biffé illisible}
\[
\struck{\pi\langle x, yF\rangle =}\ \langle Fx, y \rangle = \pi \langle x, yF \rangle ,
\qquad \langle x, yV \rangle = \pi \langle Vx, y \rangle
\]

\page{8}

\note{page entière annulée par trois longs traits presque verticaux ; elle
se lit néanmoins, et on la donne}

\emph{\uncertain{Question}} Comment définir directement une connexion
\uncertain{intégrable} sur $\mathbb{R}^{*} f_{*}(\Omega^{\bullet}_{X/S})$ ??

\emph{Lemme} $X$ topos, $K^{\bullet}$, $L^{\bullet}$ complexes abéliens sur
$X$, $\mathcal{G}$ faisceau abélien de $X$,
\[
\mathcal{G} \xrightarrow{\ i\ } \underline{\mathrm{Hom}}^{\bullet}(K^{\bullet}, L^{\bullet})^{-1}
= \prod_i \underline{\mathrm{Hom}}(K^{i}, L^{i-1})
\]
$P$ un $\mathcal{G}$-torseur,
$P \xrightarrow{\ \theta\ } Z\bigl(\underline{\mathrm{Hom}}^{\bullet}(K^{\bullet}, L^{\bullet})^{(0)}\bigr)$

\struck{Alors \ill{} \uncertain{espace} \ill{} $P$ un homomorphisme \ill{}
\ill{} $D(X)$, \ill{} \ill{} un \ill{} de $K^{\bullet} \to$}\note{deux lignes
biffées d'un trait ondulé}
\uncertain{compatible} avec les \uncertain{opérations} de
$\underline{\mathrm{Hom}}(K^{\bullet}, L^{\bullet})^{-1}$
\[
\theta(g \cdot p) = \theta(p) + \underbrace{\bigl(d \circ i(g) + i(g) \circ d\bigr)}_{\delta(i(g))}
\]

Alors on \uncertain{peut} trouver un \uncertain{homom.} \struck{\ill{}}
\uncertain{dans} $D(X)$
\[
\varphi = \varphi_{K, L}(\mathcal{G}, P, i, \theta) : K^{\bullet} \longrightarrow L^{\bullet}
\]
tel que :
\begin{enumerate}
\item[1°)] $\varphi$ est \uncertain{fonctoriel} en $L^{\bullet}$
(\struck{\uncertain{complexe}} variable dans $C^{+}(X)$) (\ill{} \ill{}
\uncertain{sens} \uncertain{évident}) ;
\item[2°)] \struck{Si $\mathcal{G}$, $P$ \ill{} \ill{}} \struck{\ill{}}
\ill{} \ill{} \ill{} \ill{} $P \xrightarrow{u} P'$, $\mathcal{G}
\xrightarrow{u} \mathcal{G}'$, $\mathcal{G}, P, i, \theta$ \uncertain{par}
$\mathcal{G}', P', i', \theta'$, avec $i = i' u$, $\theta = u \theta'$,
\struck{\ill{}} \ill{} ;
\item[3°)] Si $P$ admet une section, alors\note{la page s'arrête ici}
\end{enumerate}

\section*{Caractères des groupes de B.T.\ à réduction semi-stable virtuelle essentielle}

\note{inscrit à l'encre en haut à droite du feuillet de garde rose (page 9),
qui ne porte rien d'autre et ne reçoit pas de numéro de page}

\page{10}

\textbf{Caractère d'un groupe de B.T. sur un corps local}.

\note{les pages 10 à 13 sont un tapuscrit, tapé en bleu, sans pagination
de frappe, et non identifié comme publié ; on le donne en entier. Le titre
est de sa main, à l'encre, souligné ; ses
autres interventions à l'encre sont signalées au fil du texte. Les
surfrappes de la dactylographie (« xxx ») ne sont pas relevées}

Soit $V$ un anneau de val.\ discrète \marginal{à corps résiduel
parfait}\note{ajout à l'encre, entouré, relié par un trait à « discrète »}
; pour ce qu'on va faire, on peut le supposer complet.
\struck{à corps résiduel algébriquement clos \ill{} semble l'hypothèse
corps résiduel parfait ; on peut alors, au besoin,} (quitte :
\struck{passer au complété de l'hensélisé strict).}\note{le passage est
barré à l'encre ; « résiduel algébriquement clos » est en outre souligné
d'un tiret discontinu, et un trait brisé encadre la dernière ligne, avec
« (quitte : » écrit en marge gauche : sans doute la parenthèse « quitte à
passer au complété de l'hensélisé strict » est-elle rétablie, mais ce n'est
pas sûr. Les derniers mots de la première ligne barrée sont surchargés à
l'encre et ne se lisent pas}
On suppose que le corps résiduel $k$ de $V$ est de car.\ $p > 0$ ; le corps
des fractions $K$ de $V$ est donc de car $p$ ou $0$. On s'intéresse aux
\add{$p$-}groupes de Barsotti-Tate sur $K$, notés $M =
(M(\nu))_{\nu \geqslant 0}$.\note{« $p$- » est ajouté à la machine au-dessus
de la ligne ; les deux $\nu$ sont portés à l'encre dans les blancs de la
frappe} On dit qu'un tel groupe \emph{satisfait à la condition de
monodromie} s'il admet une suite de composition par des groupes de B.T.,
telle que les $N = \operatorname{Gr}^{i}(M)$ aient \emph{essentiellement
bonne réduction}, i.e.\ qu'il existe pour chacun d'eux une extension finie
\add{séparable} $K'$ de $K$ telle que $N_{K'}$ ait bonne réduction sur le
normalisé $S'$ de $S$ dans $K'$.\note{« $N =$ » est ajouté à la machine
au-dessus de « les $\operatorname{Gr}^{i}(M)$ » et relié par un trait à
l'encre ; « séparable » est tapé au-dessus de la ligne, entouré à l'encre}
Alors $N_{K'}$ provient donc, d'après le th.\ de Tate (démontré pour
l'instant seulement dans le cas d'\add{in}égales caractéristiques), d'un
groupe de B.T. $N$ sur $S'$ déterminé \emph{à isomorphisme unique près}.
Par suite, si $K'$ est prise galoisienne de groupe $G$, $G$ opère sur $N$ de
façon compatible avec son action sur $S'$. --- Bien entendu, on peut
prendre $K'$ indépendant de $i$ ; on dira, si $K'$ marche pour $M$
satisfaisant la condition de monodromie, que $K'$ est un \emph{corps de
bonne réduction virtuelle} pour $M$.\note{« in » est tapé au-dessus de
« d'égales » et relié à l'encre ; $S$ n'a pas été introduit : c'est
évidemment $\operatorname{Spec} V$}

Soit $K_{BT}(K; K')$ le groupe « de Grothendieck » défini par les groupes
de BT sur $K$ admettant $K'$ comme corps de bonne réduction virtuelle.
Alors, si $K'/K$ est galoisienne de groupe de Galois $G$, on trouve un
homomorphisme naturel
\[
K_{BT}(K; K') \longrightarrow K_{BT}(S'; G) , \tag{1}
\]
où le deuxième membre est le groupe de Grothendieck des groupes de B.T.
sur $S'$ avec opérations de $G$ (compatibles aux opérations de $G$ sur
$S'$).\note{le numéro (1) est à l'encre}
D'ailleurs, $S$ étant complet donc hensélien, $S'$ est local, de sorte que,

\page{11}

si $s'$ est son point fermé, on trouve aussi un homomorphisme de
restriction
\[
K_{BT}(S', G) \to K_{BT}(s', G) , \tag{2}
\]
en notant que $G$ opère sur $s'$ (de façon pas nécessairement fidèle).
Or, $k' = k(s')$ étant supposé \struck{complet} parfait, on a un foncteur
\emph{covariant} naturel
\[
\mathrm{BT}(s') \to \mathrm{Modlib}(W(k')) ,
\]
\marginal{et \uncertain{même} vers $\mathrm{Cris}(k')$}\note{« parfait » est
tapé au-dessus de « complet » surfrappé ; l'ajout « et même vers
$\mathrm{Cris}(k')$ » est au crayon, à droite de la formule}
et comme $G$ opère sur $k'$, donc sur $W(k')$, le foncteur précédent induit
trivialement un foncteur
\[
\mathrm{BT}(s', G) \to \mathrm{Modlib}(W(k'), G) ,
\]
d'où par suite un homomorphisme
\[
K_{BT}(s', G) \longrightarrow K(W(k'), G) . \tag{3}
\]
Composant (1), (2) et (3), on trouve un homomorphisme canonique
\[
K_{BT}(K; K') \longrightarrow K(W(k'), G) . \tag{4}
\]
Notons d'ailleurs que si $k$ est algébriquement clos, alors $G$ opère
trivialement sur $k'$ donc sur $W(k')$, et on a alors
\[
K(W(k'), G) = R_{W(k')}(G) ;
\]
dans le cas général, il y a lieu d'introduire le groupe
d'\uncertain{intertie}\note{sic, pour « inertie »} $I \subset G$, et on
trouve un homomorphisme
\[
K(W(k'), G) \to K(W(k'), I) \simeq R_{W(k')}(I)^{G/I} ,
\]
\marginal{\textbf{NB} $G/I$ opère sur $R$ via ses opérations sur $I$
\emph{et} sur $W(k')$ par transport de structure}
\note{l'exposant $G/I$ est ajouté à l'encre, et répété entouré au-dessus
du signe $\simeq$ ; la note NB est à droite, séparée par un trait
vertical. Dans la marge gauche, en oblique, trois lignes de crayon :
« On \uncertain{peut} \uncertain{dire} \ill{} \uncertain{suivant} les
\uncertain{définitions} \ldots »}
d'où grâce à (4)
\[
K_{BT}(K, K') \to R_{W(k')}(I)^{G/I} \longrightarrow
\operatorname{Cent}(I, W(k'))^{G} , \tag{5}
\]
\note{les exposants $G/I$ et $G$ sont à l'encre ; la seconde flèche est
redessinée à l'encre et le $k'$ du dernier terme repassé}
la dernière flèche étant la flèche trace. L'homomorphisme composé est
appelé encore l'homomorphisme trace, l'image de $\mathrm{c}\ell(M)$ dans le
dernier membre est appelé le \emph{caractère} du groupe de B.T. $M$. On
voit facilement que dans un sens évident, il ne dépend pas du choix du corps
de bonne réduction virtuelle $K'$ : la fonction caractère est alors
envisagée comme fonction sur le \add{sous-}groupe d'inertie (infini) de
$\operatorname{Gal}(\bar K/K)$.\note{« $\ell$ » de « $\mathrm{c}\ell$ » et
la barre de $\bar K$ sont portés à la main ; « sous- » est tapé au-dessus
de la ligne}

Supposons que $M_{K'}$ soit \emph{constant}, de sorte que $T_p(M)$ est un
faisceau $p$-adique localement constant (au sens technique

\page{12}

du terme). On voit alors facilement que le caractère de $M$ coïncide, sur
les éléments de $I$ \emph{d'ordre premier à} $p$, avec le \emph{caractère
modulaire} de la représentation de $I$ sur $\mathbb{F}_p$ définie par
$M(0)$ ?\note{le $\mathbb{F}$ de $\mathbb{F}_p$, le $0$ de $M(0)$ et le
point d'interrogation sont repassés à l'encre}
Comme, pour $K$ de caractéristique nulle, la donnée d'un groupe de BT sur
$K$ revient à la donnée d'une représentation $p$-adique de
$\operatorname{Gal}(\bar K/K)$ (dans un Module libre de type fini sur
$\mathbb{Z}_p$), on peut se demander s'il n'est pas vrai en tous cas que le
caractère d'un $M$ (satisfaisant la condition de monodromie) est déterminé
ainsi par la seule connaissance de $M(0) = {}_{p}M$. Il n'en est rien, car
si cette condition était vérifiée, on en déduirait facilement, en utilisant
la théorie de Brauer des relèvements virtuels des représentations sur
$\mathbb{F}_p$ en représentations sur $\mathbb{Z}_p$, que pour \emph{tout}
$M$ satisfaisant la condition de monodromie, le caractère de $M$ coïncide
avec le caractère modulaire de $M(0)$ sur les éléments d'ordre premier à
$p$ de $I$. En particulier, si $M$ a bonne réduction, il s'ensuivrait que
le caractère modulaire de $M(0)$ serait trivial dès que $M$ a bonne
réduction.\note{« de $M(0)$ » est tapé au-dessus de « serait » et relié à
l'encre}
Or c'est déjà faux en prenant $M = \mathbb{Z}_p(1)$, donc $M(0) = \mu_p$,
en se plaçant dans le cas où $K$ ne contient pas les racines $p$.èmes de
$1$.\note{le $\mu$ de $\mu_p$ est porté à l'encre dans le blanc de la
frappe}
Dans le cas où il les contient, au contraire, il n'y a pas d'ennuis pour
les groupes de BT sur $K$ qui sur un $K'$ ont une bonne réduction qui est,
soit de type constant (cf plus haut), soit de type multiplicatif, et il
conviendrait alors d'examiner encore le cas où il est de type
ind-unipotent sur la \struck{fibre} point fermé.

\marginal{\textbf{NB} Cela impliquerait que la fonction caractère
\uncertain{est} à valeurs dans $\mathbb{Z}_p \subset W(k)$, dès que $K$
contient les racines $p$.\uncertain{ièmes} de $1$ \ill{} \ill{} [\ill{}
\ill{} \uncertain{congruence} \ill{} $(p-1)$] \struck{\ill{}}}\note{note au
crayon dans la marge gauche, en face des lignes suivantes ; la dernière
ligne est biffée}

Cas test : soient $E_i$ $(i = 1, 2)$ deux courbes elliptiques sur $K$, qui
ont bonne réduction essentielle, ${}_{p}E_1$ isomorphe à ${}_{p}E_2$, $K$
contient les racines $p$.èmes de $1$, est-il vrai que les fonctions traces
$\phi_{E_1}$ et $\phi_{E_2}$ sur le groupe d'inertie (qu'on peut, comme on
sait, calculer aussi via les $T_\ell$ pour $\ell$ nombre premier
\emph{quelconque}) sont égales \add{sur les éléments premiers à $p$} ?
\note{« $(i = 1, 2)$ » est ajouté à l'encre au-dessus de « $E_1$, $E_2$ »
surfrappés ; les $\ell$ sont portés à la main ; « sur les éléments premiers
à $p$ » est tapé au-dessus de la ligne, entouré à l'encre et renvoyé après
« égales »}
Le cas intéressant est celui où l'une au moins des deux courbes, après
extension $K \to K'$, se réduit en une courbe à invariant de Hasse nul.
Ou encore : si $A_K$ schéma abélien sur $K$

\page{13}

(par exemple courbe elliptique) est tel que ${}_{p}A_K$ soit un groupe
constant, alors est-il vrai que la trace $\phi_E$ soit \add{égale à $n =
\dim A_K$} sur les \struck{points} éléments premiers à $p$ de $I$ ? (Seul le
cas de réduction non semi-stable de $A$ offre un problème ; donc il
faudrait ici trouver un exemple où le critère de réduction semi-stable de
Raynaud est en défaut pour $\ell = p$ !).\note{« égale à $n = \dim A_K$ »
est tapé au-dessus de la ligne, sur des mots surfrappés ; le $\ell$ est
porté à la main}

\marginal{faux par Tate}\note{au crayon, en oblique dans la marge gauche,
en face du paragraphe qui suit}
Dans le cas d'une courbe elliptique, et en car.\ rés.\ $\neq 2, 3$, prenant
pour $K'$ la plus petite extension sur laquelle on a réduction
semi-stable, on trouve que $G$ est un sous-groupe de $\operatorname{Aut}(A_{s'})$,
lequel \add{est cyclique} \struck{a} (2, 4 ou 6 éléments) donc est premier à
$p$, et la question revient à celle de savoir si $\phi_E$ est de valeur
constante égale à $2$ \uncertain{(?)} \struck{si} \add{lorsque} ${}_{p}A$
est constant (et si $K$ contient les racines $p$.èmes de l'unité, il suffit
de regarder le cas où la courbe elliptique réduite, après extension $K \to
K'$, est d'invariant de Hasse nul).\note{« est cyclique » est tapé
au-dessus de la ligne et relié à l'encre ; « lorsque » est à l'encre
au-dessus de « si » surfrappé, et le « (?) » à l'encre ; la parenthèse
ouvrante devant « 2, 4 ou 6 » et la parenthèse fermante finale sont
portées à l'encre}

\section*{Détermination des « stratifications canoniques » sur les déformations}

\note{inscrit au crayon en haut à droite du feuillet de garde (page 14), sous
un titre biffé de quatre traits, « Déformations des groupes de B.T.
ordinaires » ; en dessous, « L'isom. de nature globale », souligné, et
encadré : $\det T_p(A)_{\text{ét}} \simeq \det T_p(A^*)_{\text{ét}}$. Le
feuillet ne reçoit pas de numéro de page}

\page{15}

$A_0$ VA sur un corps $k$ de car.\ $> 0$.

\emph{Question} : quand est-il possible, pour \uncertain{toute} déf.\
inf.\ $A/\Lambda$ de $A_0$ (avec $\Lambda$ $k$-algèbre artinienne locale)
d'associer une $k$-stratification sur $\mathrm{DR}^1(A/\Lambda)$,
compatible avec \uncertain{changement} de base et fonctorielle en $A$ ?

Si $M$ est \uncertain{le} \uncertain{module} des \uncertain{modules}
formels \struck{\ill{}} de $A_0$ sur $k$, \struck{il faut} et $A_M$ la
déformation universelle, il faut que $\mathcal{H} = \mathrm{DR}^1(A_M/M)$
ait une \struck{\ill{}} \emph{stratification} formelle \emph{compatible
avec les opérations} $F, V$ :
\[
\mathcal{H}^{(p)} \overset{F}{\underset{V}{\rightleftarrows}} \mathcal{H} .
\]
Cela implique que pour $x \in M$, le type (\uncertain{géométrique}) de
$\mathcal{H}_x$ muni de $F_x$, $V_x$ est le même que celui de $\mathcal{H}_0
= \mathrm{DR}^1(A_0)$, muni de $F_0$, $V_0$. Cela implique sans doute que
$A_0$ est « ordinaire » i.e.

\[
\begin{array}{ccccccc}
0 \to & H^{10} & \to & H & \to & H^{01} & \to 0 \\
 & & & {\scriptstyle V}\downarrow\;\uparrow{\scriptstyle F} & & & \\
0 \to & H^{10(p)} & \to & H^{(p)} & \to & H^{01(p)} & \to 0
\end{array}
\]
\[
\operatorname{Im} V = \operatorname{Ker} F = H^{10(p)} , \qquad
\operatorname{Ker} V = \operatorname{Im} F ,
\]
\note{à gauche du trait vertical qui les sépare de la prose ; sous $V$ et
à droite de $F$, deux signes biffés illisibles. Les exposants $10$, $01$
sont lus sur des traits rapides}
\marginal{à droite du diagramme : \ill{} \uncertain{la} \uncertain{sub}
\uncertain{max.} \ill{} \uncertain{pts} \ill{} \uncertain{ordre} $p$,
\ill{} \ill{} \ill{} \ill{} \ill{}, que \uncertain{les}
$\underline{t}_{A_0^*} = H^1(A_0, \mathcal{O}_{A_0}) = H^{01}(A_0)$, la
\uncertain{puissance} $p$.\uncertain{ième}}

\uncertain{soit} \uncertain{un} \uncertain{supplémentaire} [i.e.\
\uncertain{induise} \uncertain{un} \uncertain{isom.} de $H^{01(p)}$ sur
$H^{01}$].\note{deux traits obliques séparent ce qui précède de ce qui
suit}
Mais dans \struck{\ill{}} \uncertain{le} \uncertain{cas} $A$ sur
$\Lambda$, \uncertain{dans} $\mathrm{DR}^1(A)$. Cela implique que
l'ext.\ $H = \mathrm{DR}^1(A)$ de $H^{01}$ par $H^{10}$ \uncertain{est}
\uncertain{canoniquement} \ill{}, \uncertain{par} $\operatorname{Im} F =
F(\mathcal{H}^{(p)})$, \ill{} \struck{\ill{}} $\operatorname{Im} F
\xrightarrow{\ \sim\ } H^{01}$ \uncertain{évidemment} ?
\note{« $= H^{01}$ » est ajouté au-dessus de la flèche, relié par un trait
courbe ; les trois dernières lignes sont d'une prose rapide dont seuls les
termes mathématiques sont sûrs}

\page{16}

stable par la stratification canonique. De même, comme $VF = 0$,
(\ill{} $= \operatorname{Ker} V = \operatorname{Im} F = H^{1}$)
$V$ est nul sur $H^{1}$, et se \uncertain{factorise} donc à travers
\[
H^{10(p)} = \operatorname{Ker} F ,
\]
\struck{\ill{} $FV = 0$}, et est égal à un homomorphisme \ill{}, \ill{}
que $H^{10(p)}$ est stable par la \struck{\ill{}} stratification, donc
\ill{}\note{la parenthèse de la deuxième ligne est ajoutée au-dessus de la
ligne et reliée par une accolade ; son premier terme ne se lit pas, et
« $H^{1}$ », ici et à la ligne suivante, est lu sur un « $H$ » suivi d'un
prime ou d'un $1$ (peut-être $H^{01}$). Sous « $H^{10(p)} =
\operatorname{Ker} F$ », un « $\S$ » et un mot biffés}

\emph{Attention}, cela \emph{n'implique pas} \emph{a priori} que $H^{10}$
le \emph{soit également}. Tout ce qu'on peut dire (\uncertain{admettons}
aussi pour $\Lambda$ une $\Lambda$-algèbre locale \uncertain{noeth.}\
\uncertain{complète}, corps résiduel $k$\ldots) c'est que si $S =
\operatorname{Spf}(\Lambda)$, et si on considère les \uncertain{deux}
\uncertain{images} \uncertain{inverses} $E_1$, $E_2$ de $H$ sur
\struck{\ill{}} $T = \operatorname{Spf}(\Lambda \hat\otimes \Lambda) = S
\times S$, \struck{\ill{}} et les \uncertain{deux} \uncertain{images}
\uncertain{inverses} $F_1$, $F_2$ de $H^{10}$, et considérant
l'isom.\ de stratification
\[
u : E_1 \xrightarrow{\ \sim\ } E_2 , \qquad \text{d'où} \qquad
u^{(p)} : E_1^{(p)} \to E_2^{(p)} ,
\]
alors
\[
u^{(p)}(F_1^{(p)}) = \struck{\ill{}}\, F_2^{(p)} , \quad \text{i.e.} \quad
u(F_1)^{(p)} = F_2^{(p)} .
\]
Si donc $u(F_1) = G$, $F_2 = G'$, $E = E_2$, \uncertain{on} \uncertain{a}
$G^{(p)} = G'^{(p)}$ comme \struck{\ill{}} \add{sections} de
$\operatorname{Grass}(E)$ sur $T$, on ne peut \uncertain{conclure} \ill{}
$G = G'$, \uncertain{ces} sections devenant égales par changement de
\uncertain{base}
\[
T' \xrightarrow{\ \pi\ } T , \qquad T' = T , \ \pi = \text{\uncertain{Frobenius} \uncertain{absolu}.}
\]
\note{sous $T'$, un signe d'égalité vertical et « $T$ », avec deux mots
superposés, lus « Frobenius absolu » sans certitude}
On \uncertain{veut} \uncertain{prouver} \ill{} \ill{} \ill{} \ill{}
\uncertain{conclure} $G = G'$, \uncertain{car} \uncertain{en}
\uncertain{général} $\pi$ \uncertain{n'est} \uncertain{pas} \uncertain{un}
\uncertain{épim.} Mais si \add{$S$ donc} $T$ est formellement lisse sur
$k$,\note{« $S$ donc » est ajouté au-dessus de la ligne, souligné d'un
trait courbe ; la page s'arrête sur cette virgule}

\page{17}

cela marchera (\uncertain{c'est} \uncertain{loc.}\ \uncertain{fid.}\
\uncertain{plat}). On sait \uncertain{que} \uncertain{c'est}
\uncertain{suffisant}, si on \uncertain{remonte} au cas $S = M$ (variété
des modules formels). Donc $H^{10}$ \emph{est stable} \emph{par la
stratification}. Donc $H$ admet une décomposition fonctorielle et stable
par ext.\ de la base
\[
H \simeq H^{10} + H^{01}
\]
et cette décomposition est \emph{stable par stratification}.

Donc le pb est maintenant de trouver une stratification \struck{\ill{}}
(\struck{fonctorielle} \uncertain{fonctorielle} \uncertain{dans}
\uncertain{les} \uncertain{deux} \uncertain{sens}) sur les deux morceaux
\[
H^{10} = \underline{t}_A^{\vee} \quad \text{et} \quad H^{01} = \underline{t}_{A^*} ,
\]
Les deuxièmes \add{pour $A_0$,} \uncertain{s'identifiant} au premier pour
$A_0^*$.\note{« pour $A_0$, » est ajouté au-dessus de la ligne, relié par
un trait courbe ; la phrase est d'une lecture incertaine au-delà des
formules}
Donc il faut trouver une stratification fonctorielle sur
$\underline{t}_A$. Or il y \uncertain{en} a \uncertain{une}
\uncertain{a} \uncertain{priori} \add{\uncertain{une}}, grâce à la
structure de \uncertain{puissances} \uncertain{divisées}, qui
\marginal{à préciser \ldots}
\uncertain{permet} d'\uncertain{introduire} la \uncertain{structure}
\struck{\ill{}} \uncertain{(pour $A$, $A$ \ill{})} \uncertain{naturelle}
\uncertain{disons} \uncertain{de} $\mathrm{Gl}(g)_k$. Est-ce la seule,
\struck{\ill{}} \ill{} \uncertain{sait} \uncertain{définir}
\uncertain{directement} en utilisant la \uncertain{compatibilité} avec
les \uncertain{puissances} $p$.ièmes (pour $\Lambda$, $A$ sur $X$,
\uncertain{donnés}).\note{« à préciser \ldots » est en marge gauche,
séparé du texte par un trait vertical le long des trois lignes ; dans la
parenthèse biffée, seuls les deux « $A$ » se lisent}

\textbf{Remarque}. On sait donc que la \uncertain{position} de la
\uncertain{filtration} \uncertain{par} \ill{} \ill{} \ill{} \ill{}
\uncertain{stratification} \uncertain{sur} \uncertain{la}
\uncertain{base} \add{\uncertain{\ldots} \uncertain{des}
\uncertain{déf.}}\ \uncertain{doit} \uncertain{donner} \uncertain{un}
\uncertain{principe} \uncertain{de} \uncertain{classification}
(\uncertain{infinitésimale} \ill{} \uncertain{en} \uncertain{car.}\ $p
> 0$ !\note{la remarque est marquée d'un double trait vertical dans la
marge gauche ; une insertion au-dessus de la ligne, reliée par un trait
courbe, ne se lit qu'en partie ; la parenthèse n'est pas fermée}

\page{18}

La décomposition envisagée (can., fonct.)
\[
\mathrm{DR}^1(A/S) \simeq \underline{t}_A^{\vee} + \underline{t}_{A^*}
\]
vaut pour le préschéma abélien \emph{ordinaire} sur un $S$ de car.\ $p >
0$, en même temps que les expressions
\[
\underline{t}_A \simeq \mathcal{L}_A \otimes_{\mathbb{F}_p} \underline{\mathcal{O}}_S ,
\qquad
\underline{t}_{A^*} \simeq \mathcal{L}_{A^*} \otimes_{\mathbb{F}_p} \underline{\mathcal{O}}_S
\]
où $\mathcal{L}_A$, $\mathcal{L}_{A^*}$ sont les faisceaux
\add{\uncertain{étales}} \uncertain{loc.}\ \struck{\ill{}} libres de rang
$g$ ($=$ dim.\ relative de $A/S$) des \uncertain{vecteurs}
\uncertain{invariants} sur $\mathbb{F}_p$. Si $S$ est le spectre d'un
corps parfait,\note{la lettre cursive que l'on rend $\mathcal{L}$ est
un L bouclé, lu tel ; « étales » est ajouté au-dessus de la ligne et
reporté par un trait}
\[
\mathcal{L}_{A^*} = (\underline{t}_{A^*})^{F} = H^1(A, \underline{\mathcal{O}}_A)^{F}
= H^1(A, \mathbb{Z}/p\mathbb{Z})
= \operatorname{Hom}({}_{p}A, \mathbb{Z}/p\mathbb{Z}) = \struck{\ill{}}\ {}_{p}A(k)^{\vee}
\]
\note{le $A$ de $\operatorname{Hom}({}_pA, \ldots)$ porte un petit signe
en exposant, peut-être un prime, qui n'est pas repris ensuite}
\[
\left\{
\begin{aligned}
\mathcal{L}_{A^*}(k) &= {}_{p}A(k)^{\vee} \\
\mathcal{L}_{A}(k) &= {}_{p}A^*(k)^{\vee}
\end{aligned}
\right.
\]
Donc
\[
\left\{
\begin{aligned}
\underline{t}_A^{\vee} &= \mathcal{L}_A^{\vee} \otimes_{\mathbb{F}_p} k = {}_{p}A^*(k) \otimes_{\mathbb{F}_p} k \\
\underline{t}_{A^*} &= {}_{p}A(k)^{\vee} \otimes_{\mathbb{F}_p} k
\end{aligned}
\right.
\]
\note{après $\mathcal{L}_A$ dans la première ligne, un indice surchargé
et illisible ; le dernier $k$ est écrit en capitale}

On a des relations analogues pour $S$ (de car.\ $p > 0$) quelconque :
l'hyp.\ sur $A$ implique (\uncertain{cas} \ill{}) que
\[
{}_{p}A \to S
\]
s'insère de façon unique comme \ill{}
\[
{}_{p}A \to ({}_{p}A)_{\text{ét}} \to S
\]
les deux morphismes finis localement libres, le premier radiciel surj., le
2\textsuperscript{ème} \emph{étale}.

\page{19}

Ceci dit, on a des isom.\ can.
\[
\left\{
\begin{aligned}
\mathcal{L}_A &= ({}_{p}A^*)_{\text{ét}}^{\vee} \\
\mathcal{L}_{A^*} &= ({}_{p}A)_{\text{ét}}^{\vee}
\end{aligned}
\right.
\]

\textbf{N.B.} Soit $\varphi$ une polarisation de $A$ (\uncertain{i.e.}\
une section de $\underline{\mathrm{NS}}_{A/S}$). Alors \ill{} \struck{\ill{}}
\struck{\ill{}} $\varphi$ \uncertain{définit} une section
\struck{\ill{} \ill{} $A$ \ill{} \ill{}} \struck{\ill{}} de
$\bigwedge^2 \mathrm{DR}^1(A/S)$, \struck{\ill{}}
\struck{Cela \ill{} \ill{} \ill{}}
\note{deux lignes et demie biffées, dont un « $A$ » et « de $S$ » ; le
passage qui suit, de « donc » à « Cet homomorphisme », est en outre
traversé de trois longues diagonales, sans doute une annulation}
donc \struck{\ill{}} \ill{} \ill{} \uncertain{composantes}
\[
\bigwedge^2 \mathrm{DR}^1(A/S) \simeq \bigwedge^2 \underline{t}_A^{\vee}
+ \underline{t}_A^{\vee} \otimes \underline{t}_{A^*} + \bigwedge^2 \underline{t}_{A^*}
\]
\uncertain{sur} $\underline{H}^{10} = \underline{t}_A^{\vee}$ et
$H^{01} = \underline{t}_{A^*}$ \uncertain{comme} \uncertain{facteurs}
\uncertain{isotropes}, et définit donc un \uncertain{homomorphisme}
\[
\underline{t}_A^{\vee} \otimes \underline{t}_{A^*} \longrightarrow \underline{\mathcal{O}}_S
\quad \text{i.e.} \quad \underline{t}_{A^*} \longrightarrow \underline{t}_A .
\]
Cet homomorphisme n'est \struck{\ill{} \ill{} \ill{} \ill{}} \uncertain{autre}
que l'\uncertain{homomorphisme} tangent à $\tilde\varphi : A \to A^*$
\[
\underbrace{\operatorname{Hom}(\underline{t}_A, \underline{t}_{A^*})}
\]
\note{la ligne « Cet homomorphisme \ldots\ $A^*$ » est barrée d'un long
trait ; sous $\operatorname{Hom}(\underline{t}_A, \underline{t}_{A^*})$,
une accolade renvoie à la composante médiane de la décomposition
précédente}

Ceci dit, \uncertain{c'est} \uncertain{dire} \uncertain{que} la
\uncertain{composante} \uncertain{médiane}, et \uncertain{donc}
\uncertain{autre} \uncertain{que} l'\uncertain{homom.}\ tangent
$\tilde\varphi' : \underline{t}_A \to \underline{t}_{A^*}$ de
$\tilde\varphi : A \to A^*$ défini par $\varphi$, qui se déduit
\uncertain{formellement} \uncertain{\ill{}}, par tensorisation
\struck{\ill{}} $\otimes_{\mathbb{F}_p} \underline{\mathcal{O}}_S$, par
\[
\mathcal{L}_A \to \mathcal{L}_{A^*} \quad \text{\uncertain{homom.}}\ \ldots
\]
\[
({}_{p}A)_{\text{ét}} \to ({}_{p}A^*)_{\text{ét}} \quad \text{\uncertain{induit}}
\ \text{par}\ {}_{p}\tilde\varphi : {}_{p}A \to {}_{p}A^* \ldots
\]
\note{au-dessus de $\mathcal{L}_A$ et $\mathcal{L}_{A^*}$, deux renvois
« $\wr\wr$ » vers $({}_{p}A^*)_{\text{ét}}^{\vee}$ et
$({}_{p}A)_{\text{ét}}^{\vee}$, écrits entre les deux lignes ; la page
s'arrête sur les points de suspension}

\page{20}

\emph{Re N.B.} Soit toujours $g =$ dim.\ rel.\ de $A/S$, et considérons
\struck{(\uncertain{formé} \ill{} \ill{} \ill{} \ill{} \ill{} \ill{})}
\[
\omega_A = \struck{A}\, \det(\underline{t}_A)^{-1} = \bigwedge^{g} \underline{t}_A^{\vee}
\]
\struck{(\uncertain{C'est} un faisceau \uncertain{inv.}\ \uncertain{à}
$p$-\uncertain{structure} \ill{} \uncertain{dégénérée}, i.e.\ d'\ill{}
\ill{} \struck{\ill{}} \uncertain{faisceau} \uncertain{loc.}\ \uncertain{ct}
\ill{} \ill{} \ill{}\ \ill{} \ill{} \ill{} sur $\mathbb{F}_p$,
savoir $\{\det({}_{p}A^*)_{\text{ét}}\}$.}\note{la première ligne, sous
la formule, est une insertion de deux lignes serrées, barrée d'un trait ;
après $\underline{t}_A^{\vee}$, une rature d'encre. La parenthèse qui
suit est biffée par un long trait qui court jusqu'au milieu de la page ;
sa fin est lue sous réserve}

\struck{Si \ill{} \ill{} \ill{} \ill{} $M$ \ill{} \ill{} \ill{}
\ill{} \ill{}, et \uncertain{si} $U$ \uncertain{est} l'\uncertain{ouvert}
des $s \in S$ tels que $A_s$ \uncertain{soit} \uncertain{ordinaire},
\ill{} \ill{} \ill{} $A|U$, \uncertain{une} \uncertain{polarisation}
\ill{} \ill{} \ill{} \ill{} \uncertain{section} \uncertain{canonique}
$\sigma$ \uncertain{de} \ill{}. \ill{} $U$ est \ill{} dense dans
$S$, \ill{} \ill{} \uncertain{pseudo}-\struck{\ill{}} \ill{} $\sigma$
\uncertain{commune} \ill{} \ill{} \ill{}, \uncertain{en} \ill{} \ill{}
\ill{} $S$), \ill{} \ill{} $S$ \uncertain{dans}
\uncertain{fonction} \ill{}}\note{un long passage encadré et traversé de
quatre diagonales et de deux traits horizontaux : annulé. Il porte
notamment $U$, $A|U$, $M$, une section $\sigma$, un « pseudo- »
surchargé ; on n'en donne que ce qui se lit}

Mais en fait, \uncertain{on} \uncertain{peut} \uncertain{parler}
d'une section \uncertain{de} $L^{\otimes(p-1)}$ sur $S$ ($A$ arbitraire),
\uncertain{définie} \uncertain{canoniquement} \uncertain{à} l'\uncertain{aide}
de
\[
\underline{t}_A^{(p)} \xrightarrow{\ \pi\ } \underline{t}_A
\]
\uncertain{en} \uncertain{passant} \uncertain{au} dét., d'où
\[
(\det \underline{t}_A)^{\otimes p} \longrightarrow \det \underline{t}_A
\]
d'où
\[
\sigma_A \in \Gamma(\det \underline{t}_A)^{1-p} = \Gamma\, \omega_A^{(p-1)}
\]
\note{le lot~2 (page~21) reprend exactement ici : « Cette section mérite
le nom d'invariant de Hasse de $A$ »}

\end{document}
